% ============================================================================
% Chapter 18 --- Wave 0: Foundations
% ----------------------------------------------------------------------------
% Purpose : install governance and decision authority; adopt the sovereignty
%           grid as an institutional instrument; build baseline inventory.
%           No technical change in this wave. First wave chapter --- establishes
%           the wave chapter template that Ch. 19-27 will mirror.
% Page target: ~4 pages.
% Dependencies: Ch. 3 (sovereignty grid adopted), Ch. 4 (self-assessment
%               typically conducted within this wave), Ch. 17 (vocabulary).
% ----------------------------------------------------------------------------
% Decisions registered in _coherence-EN.md:
%   D18.1 Wave 0 is the entry point for every trajectory
%   D18.2 wave chapter template established for Ch. 19-27 to mirror
%   D18.3 Wave 0 has no technical conformance subset; exit is procedural
%   D18.4 Wave 0 trajectories barely diverge --- three short subsections
% ============================================================================

\chapter{Wave~0 --- Foundations}
\label{ch:wave0-foundations}

\begin{caveatbox}[Dependencies]
\noindent Wave~0 is the entry point of the playbook. It has no
technical pre-requisite. It draws on the sovereignty grid of
\trchap{ch:sovereignty-grid}{}, which the wave formally adopts, and
on the self-assessment of \trchap{ch:self-assessment}{}, which the
wave conducts as part of its baseline inventory work. The vocabulary
follows \trchap{ch:playbook-reading}{}.
\end{caveatbox}

Wave~0 is the wave that produces no technical change. Its purpose
is to install the governance, the decision authority, and the
baseline inventory without which the technical waves cannot be
evaluated, prioritised, or financed. Institutions that skip Wave~0
typically rediscover its content midway through Wave~1 or Wave~2,
under conditions less conducive to deliberation.

%-----------------------------------------------------------------------------
\section{Goal and dependencies}
\label{sec:w0-goal}

The goal of Wave~0 is to establish four institutional artefacts that
the rest of the playbook will reference:

\begin{itemize}[leftmargin=1.5em,nosep]
  \item an executive mandate authorising the programme, with named
    sponsorship and a declared scope;
  \item a governance body --- either dedicated or composed of an
    existing institutional body extended for this purpose --- with
    decision rights over the programme's strategic choices;
  \item the formal adoption of the sovereignty grid of
    \trchap{ch:sovereignty-grid}{} as the institution's
    deliberation instrument for procurement, audit and renewal;
  \item a baseline inventory of the institution's current digital
    infrastructure, sufficient to populate the self-assessment of
    \trchap{ch:self-assessment}{} and to derive the trajectory and
    pace of subsequent waves.
\end{itemize}

\noindent No technical change is undertaken in Wave~0. The wave is
consequently the lowest-risk in the playbook and, simultaneously,
the most consequential: every subsequent wave inherits its
artefacts as constraints.

%-----------------------------------------------------------------------------
\section{Pre-requisites}
\label{sec:w0-prereq}

Wave~0 has no technical pre-requisite, but two institutional
conditions must hold for the wave to begin meaningfully:

\begin{itemize}[leftmargin=1.5em,nosep]
  \item the institution's leadership has accepted, in principle, that
    the programme will engage the sovereignty dimension of its
    digital infrastructure --- not as a technical refactoring but as
    a deliberate policy commitment;
  \item the institution's existing IT and security functions are
    represented in the governance body, so that the programme's
    decisions are made in awareness of operational reality rather
    than in parallel to it.
\end{itemize}

\paragraph{Modest budget.} A foundation budget --- indicatively a
small staffing envelope of a few staff-months to a low single-digit
number of full-time-equivalent staff over the wave's duration,
scaled to the institution's size ---
covers the inventory work, the governance secretariat, and the
external review or peer consultation that the institution may
choose to commission. Larger budgets at this stage are typically
unnecessary and may bias the programme toward action before
deliberation.

%-----------------------------------------------------------------------------
\section{Coexistence pattern}
\label{sec:w0-coexistence}

Wave~0 has no coexistence pattern in the technical sense: there is
no incumbent component to bridge against. The pattern is
\emph{institutional}: the new governance body and the new sovereignty
deliberation overlay existing IT decision processes without
displacing them. Procurement, change management, incident response
all continue under their current arrangements; the new artefacts
inform but do not yet bind those arrangements.

\paragraph{Avoiding parallel structures.} An anti-pattern at this
stage is the creation of a parallel governance structure that
duplicates rather than extends existing decision rights. The
mitigation is to integrate the new artefacts into existing bodies
where possible (an IT Council adds the sovereignty grid to its
agenda; a procurement committee adopts grid-based gating clauses)
rather than create competing forums.

%-----------------------------------------------------------------------------
\section{Trajectory: greenfield}
\label{sec:w0-greenfield}

For institutions in a genuinely greenfield position --- new
institutions, deliberate refoundations, or research initiatives
incubated outside an existing institution --- Wave~0 is the wave
where governance is built from scratch around the sovereignty grid.
The grid becomes the founding instrument, and the inventory is
trivial because there is little to inventory. The trajectory is
correspondingly compressed (typically under six months) but still
non-trivial: the absence of existing structures requires the
governance design to be more deliberate, not less.

%-----------------------------------------------------------------------------
\section{Trajectory: brownfield single-suite-centric}
\label{sec:w0-brownfield-ms}

For the starting point most common among institutions of the
adopting institution's type and jurisdiction, Wave~0 is principally
the work of \emph{integrating} the sovereignty grid into an
established procurement and governance practice that has historically
operated around the incumbent suite's enterprise-agreement renewal cycles (e.g.\ a Microsoft Enterprise Agreement). The
grid is examined against the institution's standard procurement
clauses; gaps are documented but not yet remedied; the governance
body's first task is to schedule the sovereignty grid's annual
review against the next procurement cycle. The inventory work in
this trajectory is dominated by the documentation of the incumbent
suite's estate components, their contractual envelopes, and their data
flows.

%-----------------------------------------------------------------------------
\section{Trajectory: brownfield hybrid}
\label{sec:w0-brownfield-hybrid}

For research-intensive institutions with mixed estates, Wave~0 is
the work of \emph{consolidating} a governance practice that often
exists in fragments across faculty IT, central IT, and research-
computing units. The sovereignty grid is the unifying instrument
that allows these fragments to align on a common deliberation
vocabulary. The inventory is more involved than in the
single-suite-centric trajectory because the components are more
diverse; this is a feature of the wave, not a delay.

%-----------------------------------------------------------------------------
\section{Reversibility criteria}
\label{sec:w0-reversibility}

Wave~0 is the most reversible wave of the playbook: the executive
mandate can be revoked, the governance body can be dissolved or
reabsorbed into existing bodies, and the artefacts produced are
informational rather than binding. The reversibility window is
indefinite; reversal at any later stage is operationally costless,
although the institutional cost of having undertaken the wave and
not pursued it is reputational.

\paragraph{What Wave~0 does \emph{not} commit to.} Wave~0 does not
commit the institution to any specific technical trajectory or pace.
The trajectory and pace are determined by the self-assessment
results during the wave, and they are themselves revisable at
every wave closure thereafter.

%-----------------------------------------------------------------------------
\section{Exit criteria}
\label{sec:w0-exit}

Wave~0 has no technical conformance subset, since it produces no
technical change. Exit is consequently \emph{procedural}: the wave
is closed when the four institutional artefacts of
\S\ref{sec:w0-goal} exist in writing, are accessible to the
relevant teams, and have been reviewed by the governance body in
its constitutive meeting.

\begin{itemize}[leftmargin=1.5em,nosep]
  \item Executive mandate: signed, scoped, with named sponsor.
  \item Governance body: convened, with declared decision rights
    and meeting cadence.
  \item Sovereignty grid: formally adopted, with the procurement
    and audit cadences declared.
  \item Baseline inventory: published, with the self-assessment
    results derived from it and recorded.
\end{itemize}

\paragraph{Onward planning.} The closure of Wave~0 produces the
inputs for the trajectory selection of
\trchap{ch:playbook-reading}{} and for the prioritisation of
Waves~1, 2 and 3, which can proceed in parallel after Wave~0
provided the institution's operational capacity allows it.

\begin{realworldbox}[Governance body precedents]
\noindent Public documentation of university IT governance bodies
exists but is uneven. Patterns reported in national
research-and-education-network community fora and in
EDUCAUSE include institutional
councils dedicated to IT strategy, joint senate--administration
committees, and federated bodies that include faculty IT
representatives. Among publicly documented examples are MIT's
Information Technology Governance Committee (ITGC), reporting to
the Provost and Executive Vice President, with related committees on
infrastructure, policy, and projects; and the central University
Information Services (UIS) at Cambridge supporting a collegiate
model that combines central and college-level IT arrangements. The
institution adopting the playbook should
identify the governance precedents that resonate with its own
constitutional framework rather than transposing wholesale.
\end{realworldbox}

%-----------------------------------------------------------------------------
\section{Regulatory anchoring}
\label{sec:w0-regulatory}

\noindent Wave~0 is the governance-foundation wave. It is the
principal anchor of the management-body-accountability obligation
of the applicable national NIS2 transposition (corresponding to
NIS2~Art.~20 on accountability for security measures and training)
and of the controller-responsibility duty of \loc{data-protection-regime}
(in the EU baseline, \gls{gdpr}~Art.~24) when read at institutional level. It
operationalises the management-system clauses of
ISO/IEC~27001:2022 (clauses 4 ``context of the organisation'', 5
``leadership'', 6 ``planning'', 7 ``support''), the COBIT~2019
governance objectives EDM01 (governance framework) and EDM03
(risk optimisation), and the NIST CSF \textsf{GOVERN} function.
The procurement-clause workstream initiated in this wave is also
the entry point for the supply-chain-security obligation of the
applicable national NIS2 transposition (corresponding to
NIS2~Art.~21(2)(d)), operationalised more fully at Wave~9. Detailed mapping in
\trchap{ch:regulatory-mapping}{} \S\ref{sec:rm-nis2},
\S\ref{sec:rm-iso}, \S\ref{sec:rm-cobit},
\S\ref{sec:rm-nist-csf}.
